% AI-written, AI cross-reviewed; no human sign-off; not a claim that RH is proved
% Atlas unit P035. Local manuscript v3, item 7.12, heading on printed PDF p. 79.
% UpperTheory.V.ZK.Up.torsStab
% paper2/inputs_one.tex:535–552. Verbatim excerpt; SHA-256 53ce4af38d5ce3250d9919ea443a2a99ae959d2e8bdddf030836bd13faa86e05.
% Original TeX uses the macros and cross-reference keys of the source manuscript.


% Start of atlas edition annotation. The following block is outside the original excerpt.
\paragraph{Status of named inputs in this reading.}
\noindent\nolinkurl{PositivityTheory.v8.KStability.kStability} --- in-repo constant (\path{Development/AutomorphicArithmetic/AutomorphicVarieties/LiteratureResults/KTheoryStability/HomologicalStability/KStability.lean:3--4}); not compared.\par
% End of atlas edition annotation.

\begin{formalstatement}[Stable torsion orders]\label{pn:P035}
For a totally complex field \(L\) and \(n'\ge n\ge 2m+4\), the torsion
in \(H_m(G_{n'},\mathbb Z)\) is finite, of order
\(\lvert H_m(G_n,\mathbb Z)_{\rm tors}\rvert\).
\end{formalstatement}\nodestatus{P035}
\begin{outline}
The case \(m=0\) follows from \(H_0(G,\mathbb Z)=\mathbb Z\).
For \(m>0\), apply integral homological stability between consecutive
ranks. Van der Kallen takes \(e_0=\max(1,\mathrm{sdim})\) and gives
surjectivity for \(n\ge2m+e_0-1\) and injectivity for \(n\ge2m+e_0\).
The ring of integers admits \(\mathrm{sdim}=1\); thus \(n\ge2m+1\)
suffices and includes all the stated ranks. The composite is an isomorphism and hence restricts to
an isomorphism of torsion subgroups. Finiteness at \(n\) follows from
the integral face-complex bound below. This is the stability argument
of Blueprint Z2.K4~\bploc{r7:review:Z2.K4}. See W.~van der Kallen, \emph{Homology stability
for linear groups}, \S2.2 (p.~273) and Theorem~4.11 (p.~289).
\cite{r21-vdk}
\end{outline}
